%!TEX root = tfg.tex
\chapter{Planes de Gestión}

\begin{quotation}[Polymath]{Benjamin Franklin (1706--1790)}
For every minute spent in organizing, an hour is earned.
\end{quotation}

\begin{abstract}
En este capítulo se concretan los principales planes de gestión aplicados al proyecto,
adaptando el enfoque del PMBOK a la escala de un TFG de desarrollo de videojuegos.
\end{abstract}

\section{Plan de gestión del alcance}
El alcance se gestiona a partir del \textit{Feature List} definido en Feature Driven
Development (FDD). Cada funcionalidad incluida en el proyecto debe satisfacer un criterio
simple: \textbf{¿Aporta valor directo al núcleo jugable o a la estabilidad del producto?}

Aquellas propuestas que no cumplan ese criterio, o que pongan en riesgo el calendario
disponible, se postergan explícitamente como mejoras futuras o trabajo de extensión fuera
del alcance del TFG.

\subsection{Estructura del alcance mediante FDD}
El alcance se organiza en dos niveles:

\begin{itemize}
    \item \textbf{Nivel 1 -- Overall Model (Modelo General):} Descomposición conceptual del
    sistema en cuatro bloques funcionales independientes: generación y estructura del mundo,
    jugador y enemigos, interfaz de usuario, y progresión y gestión de recursos. Este modelo
    proporciona una visión coherente sin descender aún a detalles de implementación.

    \item \textbf{Nivel 2 -- Feature List (Lista de Funcionalidades):} Descomposición de cada
    bloque en features manejables. Cada feature es una unidad de trabajo independiente que
    puede diseñarse, implementarse y testearse en una iteración o fraccionarse si es muy
    grande. La lista es priorizada considerando dependencias técnicas, riesgo y valor para
    el núcleo jugable.
\end{itemize}

\subsection{Control del alcance durante la ejecución}
Para evitar que el alcance se corrompiera:
\begin{itemize}
    \item \textbf{Evaluación estricta:} Toda propuesta nueva se evalúa inmediatamente
    contra el criterio establecido. Si no aporta valor directo al núcleo o a la estabilidad,
    se rechaza formalmente y se registra en un backlog de mejoras futuras.

    \item \textbf{Visibilidad del backlog:} El backlog se mantiene actualizado y visible en
    GitHub Projects, de forma que el equipo puede ver en cualquier momento qué ha sido
    aceptado, qué ha sido pospuesto y por qué.

    \item \textbf{Replanificación iterativa:} Si durante una iteración se descubre que una
    feature es más compleja de lo estimado, se puede subdividir y abocar la primera parte a
    la iteración actual, posponiendo la segunda a la siguiente. Esto mantiene el alcance
    predefinido pero con flexibilidad en la distribución temporal.

    \item \textbf{Escalada al equipo:} Cambios significativos en el alcance requieren
    consenso explícito en una reunión. No se toman en aislamiento; se discuten en reuniones
    de planificación de iteración.
\end{itemize}

\subsection{Nivel de detalle del alcance}
El alcance se define con suficiente detalle para ser verificable, pero no tanto como para
congelar decisiones de diseño. En otras palabras:

\begin{itemize}
    \item Se especifica \textbf{qué} debe hacer cada feature (su comportamiento observable).
    \item Se especifica \textbf{por qué} es importante (su contribución al producto).
    \item \textbf{No} se especifica en detalle el cómo (eso pertenece a la fase de diseño
    de cada iteración).
\end{itemize}

Esta aproximación permite que el equipo tenga claridad sobre las metas sin ser presionado
por decisiones de implementación que aún no están maduras.

\section{Plan de gestión del cronograma}
El cronograma del proyecto se divide en tres fases principales que ordenan el trabajo
desde la preparación inicial hasta el cierre final de la memoria. Cada fase establece
hitos específicos que marcan el progreso y permiten detectar desviaciones respecto a la
planificación inicial.

\subsection{Fase 1: Planificación}
En esta fase se establecen los cimientos del proyecto: definición clara de objetivos,
análisis detallado del alcance, selección y validación de herramientas técnicas,
configuración del repositorio con políticas de integración, y preparación de la
documentación base de la memoria.

Su finalidad es dejar preparado un marco de trabajo sólido y coherente antes de iniciar
la construcción del videojuego. Durante esta fase no hay trabajo técnico de
implementación, sino únicamente de arquitectura de proyecto, análisis de requisitos y
disposición de la infraestructura que soportará el desarrollo posterior.

\subsection{Fase 2: Ejecución, seguimiento y control}
Esta fase concentra el grueso del trabajo técnico y se estructuran en iteraciones
independientes. Las funcionalidades se abordan de forma cíclica, en entregas cerradas que
permiten diseñar, implementar, validar e integrar incrementos jugables sin acumular
demasiado trabajo abierto simultáneamente.

Cada iteración se divide en dos bloques diferenciados:

\subsubsection{Ejecución: Diseño e Implementación}
En este apartado se concreta el trabajo técnico de la iteración. Se estructura en tres
partes:

\begin{itemize}
    \item \textbf{Lista de Funcionalidades:} Definición explícita de qué \textit{features}
    se van a abordar en la iteración. Para cada funcionalidad se deja clara su
    contribución al núcleo jugable o a la estabilidad del producto, justificando por
    qué se eligió en este momento de la planificación. Este apartado también registra
    cualquier decisión relacionada con priorización, dependencias o cambios respecto a
    la planificación inicial.

    \item \textbf{Diseño:} Descripción de las decisiones de arquitectura, patrones y
    estructuras que se van a aplicar para resolver cada funcionalidad. El diseño no
    entra en detalles de implementación línea a línea, sino que presenta las ideas
    conceptuales clave, la separación de responsabilidades, las interfaces entre
    sistemas y las restricciones técnicas que se tuvieron en cuenta.

    \item \textbf{Implementación:} Narrativa del proceso real de construcción. Se describe
    cómo se tradujeron las decisiones de diseño en código, qué desafíos técnicos
    aparecieron, cómo se resolvieron y qué aprendizajes quedaron registrados para
    futuras iteraciones. Este apartado también recoge decisiones sobre integración con
    sistemas anteriores y validación de funcionamiento durante el desarrollo.
\end{itemize}

\subsubsection{Seguimiento y control: Resultados de la iteración}
En este apartado se documentan los resultados obtenidos y se valida el cumplimiento de
los objetivos de la iteración.

\begin{itemize}
    \item \textbf{Validación de Funcionalidades:} Confirmación de que cada feature se
    comporta según lo esperado dentro del motor, se integra correctamente con sistemas
    anteriores y no introduce regresiones en funcionalidades ya implementadas.

    \item \textbf{Estado Técnico:} Análisis del estado arquitectónico del sistema tras la
    iteración. Se describe qué partes de la arquitectura quedaron consolidadas, qué
    módulos se dejaron abiertos para cambios posteriores y cuáles son las hipótesis
    técnicas validadas.

    \item \textbf{Artefactos y Diagrama de Diseño:} Generación de un diagrama UML o
    equivalente que represente la estructura del sistema después de la iteración,
    facilitando la comunicación sobre la evolución de la arquitectura y permitiendo
    documentar cómo la solución crece de forma coherente.

    \item \textbf{Aprendizajes y Desviaciones:} Documentación de qué salió bien, qué
    requirió replanificación, y qué cambios al alcance o al cronograma fueron necesarios.
    Esta información retroalimenta la selección de funcionalidades para las iteraciones
    futuras.
\end{itemize}

De esta forma, cada iteración completa genera un incremento verificable del producto que
se puede jugar, testear e integrar, a la vez que deja registrado cómo se llegó a esa
solución y qué se aprendió en el camino.

\subsection{Fase 3: Cierre}
La fase final se dedica a consolidar la memoria, revisar la coherencia entre documentación
y producto implementado, elaborar manuales de usuario e instalación, preparar los
materiales de entrega y cerrar formalmente el proyecto como trabajo académico.

En esta fase no se abordan nuevas funcionalidades, sino que se cierra el producto tal
como quedó al final de la última iteración y se prepara la documentación de cierre.

\section{Plan de gestión de comunicaciones}
Aunque el equipo está formado por dos personas, se han definido canales y rutinas de
comunicación para mantener la coordinación continua, registrar decisiones relevantes y
asegurar que ambos miembros operan desde la misma comprensión del estado del proyecto.

\subsection{Canales de Comunicación}
La comunicación diaria se articula sobre tres ejes:

\begin{itemize}
    \item \textbf{Discord:} Canal principal para coordinación inmediata, resolución rápida
    de bloqueos técnicos, decisiones operativas del día a día y avisos de urgencia. Su
    naturaleza asincrónica permite que cada miembro pueda responder en su propio ritmo
    sin requerir disponibilidad simultánea.

    \item \textbf{GitHub:} Centraliza toda la trazabilidad formal del proyecto. Las
    incidencias técnicas, los cambios registrados en el repositorio, los comentarios sobre
    decisiones de arquitectura y la evolución del backlog quedan documentados de forma
    permanente, permitiendo reconstruir el razonamiento detrás de cada cambio.

    \item \textbf{Reuniones síncronas:} Se realizan cuando es estrictamente necesario
    resolver ambigüedades, alinearse en decisiones de diseño o revisar de forma colaborativa
    avances complejos. Estas reuniones quedan documentadas con notas y decisiones registradas
    para referencia posterior.
\end{itemize}

\subsection{Cadencia de Reuniones: Sincronización iterativa}
Aunque se pueden convocar reuniones en cualquier momento si surge necesidad, existe una
cadencia obligatoria de reuniones sincronizadas al inicio y al cierre de cada iteración:

\begin{itemize}
    \item \textbf{Reunión de cierre de iteración:} Se realiza al finalizar el ciclo de
    desarrollo de una iteración. En esta reunión se:
    \begin{itemize}
        \item Valida colectivamente que las funcionalidades implementadas cumplen los
        criterios de aceptación definidos.
        \item Documentan de forma consensuada los resultados y aprendizajes de la iteración.
        \item Se identifica qué tuvo que replanificarse, qué riesgos emergieron y cómo se
        gestionaron.
        \item Se elabora un resumen del estado técnico que sirvirá como línea de base para
        la siguiente iteración.
    \end{itemize}

    \item \textbf{Reunión de planificación de iteración:} Se realiza antes o al inicio de
    la siguiente iteración. En esta reunión se:
    \begin{itemize}
        \item Seleccionar las funcionalidades que se van a abordar en el próximo ciclo,
        considerando dependencias, estimaciones de esfuerzo y riesgos técnicos.
        \item Establecen criterios de aceptación claros para cada funcionalidad.
        \item Se distribuyen responsabilidades entre los dos miembros del equipo si es
        necesario.
        \item Se actualizan el \textit{backlog} y se deja documentada la decisión de
        planificación.
    \end{itemize}
\end{itemize}

De esta forma, la comunicación combina flexibilidad para el día a día con momentos
deliberados de sincronización donde se alinean las visiones, se consolida el aprendizaje
y se prepara el siguiente ciclo con claridad compartida.

\section{Plan de gestión de riesgos}
Los principales riesgos detectados en el proyecto se agrupan en dos categorías principales,
cada una con su estrategia de mitigación.

\subsection{Riesgos técnicos}
Los riesgos técnicos son aquellos relacionados con la viabilidad de las soluciones
elegidas, la estabilidad de la arquitectura y la integración entre sistemas complejos.

\begin{itemize}
    \item \textbf{Generación procedural:} La mayor incertidumbre reside en si el algoritmo
    de generación será capaz de producir niveles válidos en tiempo aceptable, manteniendo
    navegabilidad garantizada y variedad adecuada. \textit{Mitigación:} Se aborda en la
    iteración más temprana posible (iteración 1) con un prototipo completamente funcional,
    permitiendo validar la viabilidad antes de construir sobre esa base.

    \item \textbf{Rendimiento e integración:} Conforme se añaden más sistemas (generación,
    enemigos, interfaz, progresión), existe el riesgo de que el motor se vuelva inestable,
    que surjan problemas de concurrencia o que la arquitectura no escale adecuadamente.
    \textit{Mitigación:} Se mantiene un enfoque modular desde el arranque, se testea cada
    feature en contexto de integración, y se dedican iteraciones intermedias a estabilización
    y refactorización si es necesario.

    \item \textbf{Cambios en terceros:} Unity puede recibir actualizaciones que rompan
    compatibilidad o que afecten a las soluciones implementadas. \textit{Mitigación:}
    Se fija la versión del motor al inicio del proyecto y solo se consideran actualizaciones
    si aportan estabilidad o características críticas, nunca durante ciclos activos de
    desarrollo.
\end{itemize}

\subsection{Riesgos de plazo}
Los riesgos de plazo surgen de la complejidad inherente del videojuego y del tiempo
limitado disponible para un TFG.

\begin{itemize}
    \item \textbf{Sobre-alcance:} La naturaleza iterativa del desarrollo de videojuegos
    hace que sea fácil incorporar nuevas ideas durante el desarrollo que no estaban
    previstas originalmente. Esto puede corromper el alcance y comprometer el cierre
    del proyecto. \textit{Mitigación:} El alcance se gestiona estrictamente a partir del
    \textit{Feature List} definido en FDD. Cada propuesta nueva se evalúa contra la
    pregunta: ``¿Aporta valor directo al núcleo jugable o a la estabilidad?''. Si la
    respuesta es no, la propuesta se pospone como mejora futura.

    \item \textbf{Estimaciones optimistas:} El esfuerzo real de implementación puede
    exceder las estimaciones iniciales, especialmente en iteraciones tempranas donde
    aún hay incertidumbre arquitectónica. \textit{Mitigación:} Se deja margen de
    maniobra explícito en la planificación. La Fase 2 (Ejecución) no cubre el 100\% del
    tiempo disponible; se reserva buffer para absorber desviaciones. Además, tras cada
    iteración se ajustan las estimaciones futuras basándose en lo que realmente costó.

    \item \textbf{Bloqueos de dependencias:} Una funcionalidad puede descubrirse más
    compleja de lo esperado, o puede encontrarse un requisito técnico que impida avanzar
    sobre una base previa. \textit{Mitigación:} Las funcionalidades se seleccionan no
    solo por valor, sino también considerando dependencias y riesgo técnico. Las features
    más delicadas se abordan lo antes posible para tener tiempo de replanificación si
    surgen problemas. Si un trabajo se bloquea, se dispone de un \textit{backlog} de
    tareas alternativas de menor riesgo para mantener el progreso.
\end{itemize}

\subsection{Estrategia general de mitigación}
La estrategia global combina anticipación con flexibilidad:

\begin{itemize}
    \item \textbf{Prototipos tempranos:} Los sistemas más delicados se construyen primero
    en forma de prototipo, permitiendo descubrir problemas antes de que propaguen a
    sistemas dependientes.

    \item \textbf{Priorización continua:} En cada reunión de planificación de iteración,
    se revisan riesgos emergentes y se reajusta la prioridad del \textit{backlog} si es
    necesario.

    \item \textbf{Margen de maniobra en cierre:} La Fase 3 (Cierre) no comienza cuando se
    termina la última iteración prevista, sino que hay buffer de tiempo para absorber
    desviaciones. Si una iteración se extiende, se recupera tiempo de otras iteraciones
    menos críticas sin afectar el cierre de la documentación.

    \item \textbf{Definición clara del mínimo viable:} El núcleo jugable está bien definido.
    Si el tiempo comienza a escasear, el proyecto se cierra con lo que haya en ese núcleo,
    y las mejoras se postergan como trabajo futuro. Esto evita que se derrumbe todo por
    intentar hacerlo todo.
\end{itemize}

\section{Plan de gestión de la calidad}
La calidad del producto se asegura mediante una \textit{Definition of Done} (DoD) común a
todo el proyecto. La DoD es un checklist explícito que define cuáles son los criterios
que debe cumplir una funcionalidad para considerarse realmente cerrada. Sin una DoD clara,
es fácil confundir ``funcionalidad desarrollada'' con ``funcionalidad completada'', lo que
lleva a acumular trabajo parcialmente integrado que genera problemas posteriores.

\subsection{Definition of Done: Criterios de aceptación}
Una funcionalidad solo se considera cerrada cuando cumple \textbf{todos} los siguientes
criterios:

\begin{enumerate}
    \item \textbf{Compila correctamente en Unity.} El código no presenta errores de
    compilación y no genera warnings críticos que indiquen problemas de diseño o
    importación incorrecta de dependencias.

    \item \textbf{Se ha integrado en el contexto real del juego.} La funcionalidad no existe
    en aislamiento, sino que ha sido testeada dentro del motor ejecutando el juego completo.
    Se valida que la feature funciona como se esperaba cuando interactúa con otros sistemas
    ya implementados.

    \item \textbf{No introduce regresiones.} Se verifica que la funcionalidad no rompe el
    comportamiento de features anteriormente implementadas. Esto incluye probar que sistemas
    ya cerrados siguen siendo estables tras integrar el nuevo código.

    \item \textbf{Ha sido revisada por el otro miembro del equipo.} Toda feature requiere una
    revisión técnica cruzada por parte del otro desarrollador. En esta revisión se validan
    las decisiones de arquitectura, la limpieza del código y la coherencia con el resto del
    sistema, dejando constancia del contraste técnico realizado.

    \item \textbf{Cuenta con el diagrama o documentación técnica asociada.} Cada iteración
    genera un diagrama UML o equivalente que refleja cómo la arquitectura ha evolucionado.
    La documentación en la memoria explica qué se implementó, cómo, y qué decisiones de
    diseño justificaron esa implementación.

    \item \textbf{Ha superado validación funcional en sesiones de testeo.} Se realiza al
    menos una sesión de juego real donde se verifica que la funcionalidad se comporta de
    forma equilibrada y coherente desde el punto de vista jugable. Esto va más allá del
    testeo técnico: se valida que la experiencia es la esperada.
\end{enumerate}

\subsection{Aplicación práctica de la DoD}
La Definition of Done actúa como barrera entre iteraciones. Si una funcionalidad no cumple
la DoD, no se cierra formalmente, independientemente de cuánto avance técnico represente.
Esto previene dos problemas comunes:

\begin{itemize}
    \item \textbf{Deuda técnica acumulada:} Sin DoD, es tentador marcar features como
    ``cerradas'' cuando están 80\% funcionales con la esperanza de completarlas después.
    En la práctica, esa deuda se acumula, y conforme pasan iteraciones, cada feature
    incompleta contamina más sistemas, haciendo que el código sea cada vez más frágil.

    \item \textbf{Sorpresas en cierre:} Sin revisión y validación claras, los problemas se
    descubren tarde. La DoD asegura que cada ciclo termina con un producto verificado,
    no con una esperanza de que funcione.
\end{itemize}

\subsection{Seguimiento de la calidad}
Dado el tamaño del equipo y la escala del proyecto, el seguimiento de la calidad se plantea
de forma práctica y no mediante un sistema pesado de métricas formales. El control se apoya
en tres comprobaciones recurrentes:

\begin{itemize}
    \item \textbf{Compilación estable en Unity:} Cada iteración debe dejar el proyecto en un
    estado ejecutable y sin errores críticos de compilación.

    \item \textbf{Validación de integración:} Se comprueba que la nueva funcionalidad no
    rompe sistemas previamente cerrados y que se comporta correctamente dentro del juego
    completo.

    \item \textbf{Sesiones de prueba jugable:} Se realizan pruebas manuales dentro del
    proyecto para verificar que la experiencia resultante sigue siendo coherente desde el
    punto de vista técnico y jugable.
\end{itemize}

Este enfoque es más realista para un TFG de dos personas que un sistema de medición más
formalizado, y resulta suficiente para sostener la calidad del producto dentro del contexto
del proyecto.

\section{Plan de gestión de costes}
La estimación económica del proyecto se apoya en tres componentes: coste de personal,
amortización de equipamiento y gastos indirectos. Aunque se trata de un trabajo académico
sin transacción económica real, esta aproximación permite valorar el esfuerzo invertido
con criterios profesionales, facilitando comparativas con proyectos reales de similar
envergadura.

\subsection{Estructura de costes}
Los costes se desglosan de la siguiente manera:

\begin{itemize}
    \item \textbf{Costes de personal:} Se estima una dedicación aproximada de 300 horas por
    integrante del equipo (dos desarrolladores), lo que suma 600 horas totales. Se adopta
    un perfil de ingeniero de software junior con un salario bruto estimado según mercado
    actual. Aunque ninguno de los dos miembros recibe pago real por este trabajo, la
    estimación refleja cuál sería el coste de ejecutar el proyecto en un contexto
    profesional.

    \item \textbf{Costes sociales:} Se aplica un factor de 30\% sobre el coste bruto de
    personal, que corresponde típicamente a cotizaciones a la Seguridad Social, seguros y
    otros costes inherentes a la relación laboral.

    \item \textbf{Costes de equipamiento:} Se realiza una amortización proporcional del
    hardware utilizado (ordenadores, periféricos, etc.) distribuida sobre el período de
    duración del proyecto. Solo se contabilizan los equipos empleados directamente en el
    desarrollo.

    \item \textbf{Costes de software y servicios:} Se registran como cero, ya que se utilizan
    principalmente herramientas de código abierto o con licencias educativas (Unity,
    GitHub, etc.). Si fuera necesario pagar por servicios de terceros, se registraría aquí.

    \item \textbf{Costes indirectos:} Se añade un porcentaje de 10\% sobre los costes
    directos para contabilizar gastos asociados al entorno de trabajo, electricidad,
    conectividad a internet y otros costes indirectos del espacio de desarrollo.
\end{itemize}

\subsection{Justificación de la estimación}
La estimación de 300 horas por miembro se basa en el análisis del alcance del proyecto.
Con aproximadamente 8-10 funcionalidades principales a abordar, 8 iteraciones de desarrollo,
más fases de planificación y cierre, 300 horas representa un promedio de 37.5 horas por
semana durante diez semanas de desarrollo intensivo, lo que es coherente con la dedicación
esperada de un TFG a tiempo completo o equivalente.
